\section{The \textsf{unique-copy} Algorithm}
\label{unique-copy}

We now study the \texttt{unique-copy} algorithm from ``ACSL By Example''.
This algorithm copies the contents from a source sequence to a destination
sequence whilst removing elements consecutive groups of duplicate elements once
the first one is copied. The return value is the length of the range.

\subsection{Adapting the Axiomatics}

The axiomatization provided in ``ACSL By Example'' is:

\listingname{unique-copy.axioms}
\cinput{uniquecopy/original.axioms}

The predicate \texttt{WhitherUnique} is defined by a \texttt{read} clause.
It must be adapted for the current version of the \textsf{WP} plug-in which
only implements a limited subset of \texttt{read} clauses:

\listingname{unique-copy.axioms [adapted]}
\cinput{uniquecopy/uniquecopy.axioms}

\subsection{Proving the Algorithm}

The specification of the \texttt{unique-copy} algorithm bases itself on the
\texttt{UniqueCopy} predicate, itself based on the \texttt{WhitherUnique}
function:

\listingname{uniquecopy.h}
\cinput{uniquecopy/uniquecopy.spec}

\clearpage
The implementation of the algorithm is:

\listingname{uniquecopy.c}
\cinput{uniquecopy/uniquecopy.impl}

The implementation can be partially proved correct against its specification by
running the \textsf{WP} plug-in: 
\fclogs{mutating}{uniquecopy}
